\documentclass[11pt]{article}

\usepackage[margin=1in]{geometry}
\usepackage{times}
\usepackage{graphicx}
\usepackage{booktabs}
\usepackage{amsmath}
\usepackage{hyperref}
\usepackage{caption}
\usepackage{enumitem}
\usepackage{natbib}

\title{Drowning in Meaning: Steering Figurative-vs-Literal\\ Interpretation Beyond Idioms}
\author{Itamar Vasserman (ID 205949985)}
\date{20/08/2026}

\begin{document}
\maketitle
\vspace{-2em}
\begin{center}
Submitted as final project report for the NLP course, RU, 2026
\end{center}
\vspace{1em}

\section{Introduction}

Idioms such as \emph{``take a bow''} or \emph{``green light''} admit two readings: a
figurative, conventionalized meaning, and a literal, compositional one. \citet{idiosteer}
(\emph{Steering Idioms}, IdioSteer) recently showed that this ambiguity can be controlled
mechanistically: a single mean-difference vector, extracted from an LLM's residual stream
and added back in at generation time, reliably pushes a model's continuation of an
ambiguous idiom prefix from a figurative reading toward a literal one (and vice versa),
while largely preserving fluency.

That result raises a natural question, and the one this project addresses: \textbf{is the
direction IdioSteer discovers a general figurative-vs-literal axis in the model's
representation space, or is it specific to the closed, conventionalized class of
expressions we call idioms?} Idioms are a special case of figurative language: frozen,
non-compositional multiword expressions with a fixed conventional meaning.
Metaphors and similes are figurative in the same broad sense (a literal reading exists
alongside an intended non-literal one) but are typically novel, compositional, and
constructed on the fly rather than memorized as a unit. If the steering vector encodes a
general semantic distinction between ``interpret this compositionally'' and ``interpret
this as a fixed non-literal unit,'' we would expect it to transfer to metaphors and
similes as well. If it instead encodes something closer to ``this specific memorized
idiom has a non-literal sense,'' we would expect little or no transfer.

We test this by (1) reproducing the IdioSteer pipeline (steering vector construction,
cross-layer additive injection, and LLM-judge evaluation) on a base LLM and a
paper-scale idiom benchmark, and (2) applying the resulting, idiom-derived steering
vector, \emph{unmodified}, to a new hand-built, dual-annotator-validated benchmark of 100
non-idiomatic figurative prefixes (conventional metaphors, novel metaphors, and similes)
with the same ambiguous-prefix structure and the same paper-matching size (20
expressions $\times$ 5 variants, mirroring the idiom side's 20 idioms $\times$ 5
variants). We compare the figurative/literal rate shift induced by the same vector, at
the same coefficients, across idioms and each figurative category.

Our main finding is a \textbf{broad, largely bidirectional generalization}: the idiom
vector shifts literal rate in the expected direction, on data it was never trained on,
across the pooled figurative benchmark, with Simile transferring even more strongly than
idioms themselves, Conventional Metaphor transferring more modestly, and Novel Metaphor
showing a flat effect on the literal-pushing side (plausibly a ceiling effect, discussed
in Section~\ref{sec:analysis}). Our own idiom replication is also directionally and
qualitatively consistent with IdioSteer's reported effect (monotonic, bidirectional,
high-coherence), though its absolute baseline literal rate remains well above the paper's
own reported baseline, a discrepancy we attribute to benchmark construction rather than
judge quality (Section~\ref{sec:analysis}).

\section{Related Work}

\textbf{Steering Idioms (IdioSteer).} \citet{idiosteer} construct a steering vector as the
mean difference between residual-stream activations collected in figurative-context and
literal-context idiom sentences, at a fixed source layer, and inject it additively at an
earlier layer during generation (cross-layer steering, found to outperform same-layer
injection). They evaluate on a 20-idiom, 100-sentence ambiguous-prefix benchmark using an
LLM-as-judge (Gemma-4-31B-it, selected by benchmarking several candidates against a
280-item human-annotated gold set with $\kappa=0.867$ inter-annotator agreement; the
chosen judge reaches 90.0\% accuracy and $\kappa=0.821$ against gold) and report a sharp,
coefficient-dependent shift in literal rate. Our project reuses this exact protocol
(vector construction, injection point, coefficient grid, and judge-based evaluation) and
asks whether it generalizes beyond the idiom class it was built on.

\textbf{IdioLink.} \citet{idiolink} introduce a retrieval benchmark linking idiomatic
spans to conceptually equivalent literal paraphrases across 107 idioms, 10{,}700
documents, and 2{,}140 queries. We use IdioLink's idiom/paraphrase pairs as the source
pool for constructing our steering vector's contrastive (figurative vs. literal)
activations, rather than for retrieval, and note it as a natural conceptual bridge between
the idiom-linking and idiom-steering lines of work.

\textbf{Activation steering more broadly.} Residual-stream activation addition
\citep{actadd} and its contrastive-prompt formalization, Contrastive Activation Addition
\citep{caa}, are the constructions IdioSteer's own mean-difference vector follows, and
which we in turn reuse for our steering vector. This line of work has otherwise mostly
targeted surface-level or coarse behavioral properties (sentiment, refusal, truthfulness,
writing style), and recent work has raised the concern that not every
semantic direction is equally steerable: some target behaviors are not linearly
reachable from a fixed base direction, a phenomenon discussed under \emph{steerability} and
\emph{non-surjectivity} \citep{steerability}. This is directly relevant to our null result
on novel metaphors (Section~\ref{sec:analysis}): a flat effect in one direction is
consistent with a ceiling/floor effect in the data, but is also consistent with the
idiom-derived direction being only partially aligned with whatever direction would move
novel-metaphor readings.

\section{Methodology}

\subsection{Steering framework}

We follow IdioSteer's protocol directly. Given paired activations $h_i^{\text{fig}}$ and
$h_i^{\text{lit}}$ (residual stream at a source layer $L_s$, at the final token of the
idiom/expression span, for the same item under a figurative- vs. literal-biasing context),
we compute the mean-difference vector (following the paper's convention, oriented toward
the figurative class)
\[
v = \frac{1}{N}\sum_{i=1}^{N} \left(h_i^{\text{fig}} - h_i^{\text{lit}}\right), \qquad
v_{\text{unit}} = v / \lVert v \rVert.
\]
At generation time we intercept the residual stream at an earlier injection layer $L_i$
and add $\alpha \cdot v_{\text{unit}}$ at the final prompt token only (a single-token,
prompt-pass intervention; we do not re-inject at every generated token). Positive
$\alpha$ pushes toward the figurative class, negative $\alpha$ toward literal. Following
IdioSteer's calibrated configuration for our model, we use $L_s = 14$, $L_i = 2$
(cross-layer injection), and a coefficient grid $\alpha \in
\{-4.78,-3.59,-2.39,-1.20,0,+1.20,+2.39,+3.59,+4.78\}$, with $n=5$ generated samples per
(item, $\alpha$) cell. The vector is built once, on IdioLink data
(Section~\ref{sec:idiom-data}), and reused unmodified for every experiment in this paper,
including the figurative-benchmark evaluation: it is never retrained on figurative
data.

Model: \texttt{meta-llama/Llama-3.2-3B}, base (non-instruct) checkpoint, fp16. We use a
base rather than instruction-tuned model to match IdioSteer's setup and to keep
continuations free-form rather than shaped by chat templates. Generation uses temperature
0.7, top-$p$ 0.9, and up to 50 new tokens per completion, matching IdioSteer's reported
generation settings exactly.

\subsection{Idiom dataset}
\label{sec:idiom-data}

The steering vector's contrastive pairs are drawn from IdioLink \citep{idiolink} via its
public \texttt{Intellexus/IdioLink} release, combining the \texttt{indexes} and
\texttt{queries} configs and applying the paper's exact-surface-form filter so that each
pair genuinely contrasts a figurative and a literal use of the same idiom. This mirrors
IdioSteer's own training pool construction, which retains 1{,}683 IdioLink sentences
across 22 idioms after the same filter (a moderately imbalanced $\sim$37\% figurative /
$\sim$63\% literal split); their evaluation idioms are entirely disjoint from this
training set, and we follow the same idiom-disjoint separation between vector
construction and evaluation. For evaluation, we built a 20-idiom, ambiguous-prefix
benchmark matching IdioSteer's own benchmark size and idiom inventory exactly: 100
sentences, 5 wording-diverse variants each, over the 20 idioms IdioSteer selected from
MAGPIE \citep{magpie}: \emph{up in the air, in
hot water, out of the blue, break the ice, behind bars, on ice, see red, pull someone's
leg, at the crossroads, lie low, shut the door on, under the sun, take a bow, have a
ball, in stitches, in the pink, green light, rough edges, in the fast lane, spill the
beans} (idiom list and construction constraints extracted directly from the paper PDF;
IdioSteer's exact released stimuli were not available to us, so this is a
hand-reconstructed set in the same style). Each prefix ends at the idiom itself and is
written so that a reader cannot tell, from the prefix alone, whether a figurative or
literal continuation is coming (e.g. \emph{``After months of delays, the city council
finally gave the project the green light''}).

\subsection{Non-idiomatic figurative benchmark}
\label{sec:fig-data}

We built a 100-item benchmark of non-idiomatic figurative expressions, matching the idiom
benchmark's size and structure exactly: 20 expressions $\times$ 5 wording-diverse
ambiguous-prefix variants, split into three categories: 10 Conventional Metaphors, 5
Novel (invented) Metaphors, and 5 Similes (50 / 25 / 25 items respectively; this
allocation was fixed early in benchmark construction and is not itself a designed
variable, unlike the categories' content, see Limitations). The three
categories are a controlled comparison, not an arbitrary sample: Conventional and Novel
Metaphor hold syntactic form fixed (a bare verb-phrase assertion, like an idiom) while
varying conventionality, isolating whether transfer depends on the model having seen the
figurative sense before; Simile instead varies syntactic form, replacing the bare
assertion with an explicit \emph{like}/\emph{as} comparison, testing whether transfer
depends on that idiom-like packaging specifically. Each item
mirrors the idiom benchmark's construction constraints: an ambiguous prefix ending
exactly at the figurative expression (no post-idiom text), with both a plausible literal
continuation and a plausible figurative continuation by design (the two need not be
equally likely, only both plausible). For example:

\begin{quote}
\emph{``Even as tempers flared around the campfire, he couldn't resist the urge to feed
the fire''}. Figurative continuation: he keeps escalating the argument, egging on the
rising tempers. Literal continuation: he keeps tossing more wood onto the actual
campfire flames.
\end{quote}

Conventional metaphors reuse widely lexicalized figurative senses (\emph{feed the fire},
\emph{build a wall}, \emph{cut the ties}, \emph{plant the seeds}, \emph{steer the ship})
that are frequent and familiar but, unlike idioms, remain compositional. Following the
paper's own non-compositionality criterion for idiomaticity, we checked each of the ten
against the standard substitutability diagnostic for separating a live conventional
metaphor from a frozen idiom \citep{nunberg}: \emph{plant the seeds of doubt/hope/rebellion}
all parse compositionally, whereas a true idiom such as \emph{spill the beans} does not
survive substituting \emph{peas} for \emph{beans}; all ten Conventional Metaphor
expressions pass. Novel metaphors invent non-lexicalized figurative extensions
(\emph{overwrite the DNA}, \emph{drain the aquarium}, \emph{recalibrate the compass}) to
probe whether compositionality, not just figurativeness, matters; similes use explicit
\emph{like}/\emph{as} comparisons (\emph{like a severed cable}, \emph{like an empty
vault}, \emph{like a corrupted file}). Each item carries a \texttt{gold} flag; 20 rows
(one rigorously-validated variant per expression) form a higher-confidence subset used as
a separate, cleaner analysis slice (Section~\ref{sec:results}).

Both benchmarks (idiom and figurative) went through a dual-annotator validation protocol
adapted directly from IdioSteer's own validation methodology: two independent passes
each write a figurative and a literal continuation for every prefix and flag any row
where one reading is not genuinely plausible (full prompt in
\texttt{validation/annotator\_prompt.md}). Every row in both final benchmarks is
confirmed bidirectionally ambiguous by this protocol.

\subsection{Evaluation protocol}

We adapt IdioSteer's LLM-as-judge protocol: given a prefix and a generated continuation,
the judge is asked to (i) state the intended figurative meaning, (ii) state the literal
meaning, (iii) check coherence, and (iv) then decide a label in
\{\textsc{figurative}, \textsc{literal}, \textsc{incoherent}\}, mirroring IdioSteer's
own judge structure and its two documented judge failure modes (not conflating a literal
\emph{physical} action with a literal reading when it happens to realize the idiom's own
conventional frame, e.g. bowing on stage for \emph{take a bow}; not over-labeling
unusual-but-coherent literal continuations as incoherent).

We label continuations with a 4-bit-quantized \textsc{Qwen/Qwen3-14B} judge, the largest
local judge that reliably fits a free-tier GPU alongside the 3B generation model. It is
part of IdioSteer's own candidate judge pool (Appendix J), though not its winner. All
quantitative results in Section~\ref{sec:results} use this Qwen3-14B judge, applied to
the final, dual-annotator-validated 100-item version of both benchmarks. We report, per
(dataset, category, $\alpha$) cell: figurative rate, literal rate, incoherent rate, and
coherence rate (figurative + literal). We additionally hand-labeled a small subset of raw
idiom continuations ourselves, independently of the judge, to obtain a direct
judge-vs-human calibration number, reported alongside the baseline-gap discussion in
Section~\ref{sec:analysis}.

\section{Experiments and Results}
\label{sec:results}

\subsection{Idiom baseline replication}

Table~\ref{tab:pooled} (idiom columns) shows literal rate by coefficient, pooled over all
20 idioms $\times$ 5 variants ($n=500$ generations per $\alpha$). Literal rate is 0.436 at
$\alpha=0$ and rises to 0.576--0.618 across every negative-$\alpha$ cell, peaking at 0.618
($\alpha=-3.59$ and $\alpha=-2.39$), while every positive-$\alpha$ cell (0.386--0.436) sits
at or below baseline, showing a clean, largely monotonic, bidirectional dose-response:
negative $\alpha$ consistently pushes toward literal, positive $\alpha$ consistently
pushes toward or past the figurative baseline. Coherence stays $\geq$0.916 in every cell,
so the vector is not degrading fluency to achieve this shift. The shape of this result
(monotonic-ish, bidirectional, high-coherence) is a good qualitative match to IdioSteer's
own reported effect. The one respect in which it still diverges is absolute level: our baseline
(0.436) sits far above IdioSteer's own reported Llama-3.2-3B baseline at this exact
configuration ($\approx$0.12, rising to $\approx$0.49 at $\alpha=-4.78$), so our benchmark
starts closer to 50/50 rather than strongly figurative-skewed, leaving comparatively less
headroom for the literal-pushing direction. We discuss this gap in
Section~\ref{sec:analysis}.

\subsection{Applying the idiom vector to the figurative benchmark}

Table~\ref{tab:pooled} (figurative columns) shows the same idiom-derived vector applied,
unmodified, to the 100-item figurative benchmark ($n=500$ per $\alpha$). The pattern is
the same shape as idioms and, if anything, cleaner: every negative-$\alpha$ cell
(0.658--0.696) sits above the $\alpha=0$ baseline (0.580), and every positive-$\alpha$
cell (0.530--0.554) sits below it. The idiom-derived vector \emph{does} shift
figurative/literal rate on data it was never built from, in the direction consistent
with IdioSteer's own convention (negative $\alpha \to$ more literal), and with an effect
size (+0.116 at $\alpha=-4.78$ vs. baseline) about 70\% as large as the idiom effect
itself (+0.166). Coherence again stays high throughout ($\geq$0.904).

\begin{table}[h]
\centering
\small
\begin{tabular}{@{}rrrrr@{}}
\toprule
$\alpha$ & idiom literal rate & figurative literal rate & idiom coherence & figurative coherence \\
\midrule
$-4.78$ & 0.602 & 0.696 & 0.916 & 0.942 \\
$-3.59$ & 0.618 & 0.662 & 0.948 & 0.934 \\
$-2.39$ & 0.618 & 0.658 & 0.922 & 0.938 \\
$-1.20$ & 0.576 & 0.678 & 0.964 & 0.956 \\
$0.00$  & 0.436 & 0.580 & 0.968 & 0.946 \\
$+1.20$ & 0.436 & 0.530 & 0.952 & 0.938 \\
$+2.39$ & 0.416 & 0.554 & 0.956 & 0.926 \\
$+3.59$ & 0.386 & 0.544 & 0.950 & 0.912 \\
$+4.78$ & 0.434 & 0.546 & 0.956 & 0.904 \\
\bottomrule
\end{tabular}
\caption{Literal and coherence rate by steering coefficient, pooled over the full
100-item idiom benchmark and the full 100-item figurative benchmark ($n=500$/cell each).
Negative $\alpha$ is the literal-pushing direction.}
\label{tab:pooled}
\end{table}

\subsection{Breakdown by figurative category}

Table~\ref{tab:category} splits the figurative-benchmark result by category. \textbf{Simile}
shows the largest literal-rate shift of any dataset in this study ($+0.256$, larger than
idioms' own $+0.166$); \textbf{Conventional Metaphor} shifts in the same direction, more
modestly ($+0.108$); \textbf{Novel Metaphor} again shows a flat, slightly negative delta
($-0.008$) on the literal-pushing side. Novel metaphors also have the most
literal-skewed baseline of the three categories (0.736, vs.\ 0.548 conventional / 0.488
simile) and the highest incoherent rate at extreme coefficients (0.088 at
$\alpha\in\{+3.59,+4.78\}$, vs.\ 0.040 at $\alpha=0$), suggesting this category is closer
to a ceiling where a further literal push has little room to register, rather than being
unaffected by steering \emph{per se}. Consistent with that reading, novel metaphors do
move in the expected direction on the figurative-pushing side: literal rate falls from
0.736 at baseline to 0.512--0.688 across $\alpha \in \{+1.20,\dots,+3.59\}$.

\begin{table}[h]
\centering
\small
\begin{tabular}{@{}lrrr@{}}
\toprule
Category & baseline literal rate ($\alpha=0$) & literal rate at $\alpha=-4.78$ & $\Delta$ \\
\midrule
Simile & 0.488 & 0.744 & $+0.256$ \\
Conventional Metaphor & 0.548 & 0.656 & $+0.108$ \\
Novel Metaphor & 0.736 & 0.728 & $-0.008$ \\
\bottomrule
\end{tabular}
\caption{Per-category literal rate shift under the maximal literal-pushing coefficient.
Full per-$\alpha$ breakdown in \texttt{analysis/summary\_by\_category\_alpha.csv}.}
\label{tab:category}
\end{table}

\begin{figure}[h]
\centering
\includegraphics[width=0.75\textwidth]{figures/literal_rate_vs_alpha.png}
\caption{Literal rate vs.\ steering coefficient $\alpha$: idioms vs.\ each figurative
category, all under the single idiom-derived steering vector.}
\label{fig:main}
\end{figure}

\textbf{Bottom line per category.} Simile: steering works, transferring even more
strongly than on idioms themselves. Conventional Metaphor: steering works, transferring
in the same direction as idioms but more modestly. Novel Metaphor: steering still works,
but the literal-pushing direction is ceiling-limited by an already literal-skewed
baseline, so the effect is only clearly visible when pushing toward figurative - not
evidence that this category is unaffected by the vector.

\subsection{Gold subset: the cleanest curve in the study}

The 20-item gold subset (one rigorously-validated variant per figurative expression)
shows the single cleanest dose-response curve of the entire project (Table~\ref{tab:gold}):
figurative rate rises almost monotonically from 0.26 at $\alpha=-4.78$ to 0.45 at
$\alpha=+3.59$ (dipping slightly at the most extreme $\alpha=+4.78$), with literal rate
falling the mirror image, and coherence staying $\geq$0.90 throughout. That the
highest-confidence slice of the figurative benchmark shows the cleanest trend in the
whole study, cleaner than either the pooled figurative benchmark or the idiom
benchmark itself, is a meaningful positive signal: it is not what we would expect if
the idiom vector's apparent transfer to figurative language were purely a benchmark or
judge artifact.

\begin{table}[h]
\centering
\small
\begin{tabular}{@{}rrrr@{}}
\toprule
$\alpha$ & figurative rate & literal rate & incoherent rate \\
\midrule
$-4.78$ & 0.26 & 0.70 & 0.04 \\
$-3.59$ & 0.26 & 0.69 & 0.05 \\
$-2.39$ & 0.26 & 0.70 & 0.04 \\
$-1.20$ & 0.25 & 0.73 & 0.02 \\
$0.00$  & 0.31 & 0.64 & 0.05 \\
$+1.20$ & 0.37 & 0.55 & 0.08 \\
$+2.39$ & 0.38 & 0.53 & 0.09 \\
$+3.59$ & 0.45 & 0.49 & 0.06 \\
$+4.78$ & 0.34 & 0.56 & 0.10 \\
\bottomrule
\end{tabular}
\caption{Gold subset (20 expressions, one variant each, $n=100$/cell).}
\label{tab:gold}
\end{table}

\section{Analysis and Discussion}
\label{sec:analysis}

\textbf{Qualitative flips.} At the item level (majority vote over 5 samples), 31 of 100
idiom (idiom, variant) pairs flip from majority-figurative at $\alpha=0$ to
majority-literal at $\alpha=-4.78$, spread across 15 of the 20 idioms, and 24 of 100
figurative (category, expression, variant) triples show the same pattern (12
Conventional Metaphor, 10 Simile, 2 Novel Metaphor), a category split that mirrors the
quantitative deltas in Table~\ref{tab:category} closely. A representative idiom flip,
\emph{on ice} (variant 1):

\begin{quote}
\emph{``Unsure if the deal would go through, the executives kept a bottle of champagne on
ice''}
\begin{itemize}[nosep]
\item Baseline ($\alpha=0$, judged \textsc{figurative}): \emph{``...ready to celebrate the
deal with a bottle of champagne, but they were unsure if the deal would go through...''}
\item Steered ($\alpha=-4.78$, judged \textsc{literal}): \emph{``...hooks on the wall. When
the deal was done, they opened the bottle and shared a toast. But the champagne bottle
was never returned to the wall...''}
\end{itemize}
\end{quote}

And a representative figurative flip, \emph{like a severed cable} (Simile, variant 4):

\begin{quote}
\emph{``Following the CEO's abrupt resignation, the connection between headquarters and
the regional branch felt like a severed cable''}
\begin{itemize}[nosep]
\item Baseline ($\alpha=0$): mixed across samples (\textsc{literal}, \textsc{figurative},
\textsc{incoherent}).
\item Steered ($\alpha=-4.78$, judged \textsc{literal} across all sampled continuations):
\emph{``...cord. The regional branch...was now left to fend for itself. This sense of
isolation and lack of support was exacerbated...''}
\end{itemize}
\end{quote}

Both examples show the same qualitative pattern: the steered continuation reinterprets the
expression's component words compositionally (an actual bottle on a physical shelf; an
actual snapped cable), exactly the intended literal reading. On the second example, this
is a direction the idiom-trained vector never saw during vector construction.

\textbf{Why is the idiom replication's baseline still far from the paper's?} Our idiom
benchmark, matched to the paper's full 100-item size, produces a shape (bidirectional,
high-coherence, largely monotonic) that is a good qualitative match to IdioSteer's
reported effect. The remaining gap is
purely in absolute level: our unsteered baseline (0.436) sits close to 50/50, while
IdioSteer's own baseline at this exact configuration is $\approx$0.12: their
ambiguous prefixes are considerably more figurative-skewed \emph{before any steering is
applied} than ours are. Two non-exclusive explanations are consistent with this.
First, judge differences: IdioSteer's final judge is a much larger model
(Gemma-4-31B-it) calibrated against 280 human-annotated items with $\kappa=0.821$
agreement to gold. To get an analogous number for our own judge, we hand-labeled a
subset of 86 raw idiom continuations ourselves (blind to the judge's own label, roughly
9--10 per $\alpha$) and compared against the Qwen3-14B judge's label on the same items:
raw agreement 74.4\% (64/86), Cohen's $\kappa=0.494$, well below IdioSteer's $\kappa=0.821$
on a much smaller sample. The disagreements are not symmetric: among items we labeled
\textsc{figurative}, the judge called them \textsc{literal} 37.2\% of the time (16/43),
versus only 11.6\% (5/43) the other way, a measurable \emph{literal-leaning} bias. A judge
that over-calls \textsc{literal} would mechanically inflate literal rate at every $\alpha$
including baseline, and part of the lower $\kappa$ is plausibly attributable to judge
selection and capacity rather than to our judging protocol itself: IdioSteer's
$\kappa=0.821$ reflects the single best performer out of $\sim$9 benchmarked candidates,
tuned across $\sim$14 prompt variants against their own gold set, whereas our Qwen3-14B
judge was chosen for GPU-memory feasibility, was not the winner among IdioSteer's own
candidates, is 4-bit-quantized, and was not separately prompt-tuned. Given this, judge
quality now looks like a real and possibly substantial contributor, not merely a
secondary confound. Second, our prefixes, while validated as bidirectionally plausible by
the dual-annotator protocol (run with two different models, independently of each other
and of whichever model authored the prefixes), may simply write in a less
figurative-default register than the paper's own (unreleased) stimuli. ``Bidirectionally
plausible'' does not pin down where the unsteered default sits, only that both readings
are reachable. We cannot fully disentangle these two factors with the data collected so
far.

\textbf{Does the idiom vector reflect a general axis?} The evidence leans toward yes, and
the two axes the benchmark was built to isolate (Section~\ref{sec:fig-data}) each show a
distinct, interpretable pattern rather than a uniform generalizes-or-not verdict. Along
the \emph{conventionality} axis (syntactic form held fixed as a bare assertion), transfer
strength tracks conventionality directly: Idiom ($+0.166$) $>$ Conventional Metaphor
($+0.108$) $>$ Novel Metaphor ($\approx 0$ on the literal-pushing side), consistent with
the vector partly encoding ``has an entrenched non-literal sense'' rather than a pure,
register-independent literal/figurative distinction, though Novel Metaphor's ceiling
effect (Section~\ref{sec:results}) means we cannot cleanly separate a genuinely weaker
effect from a baseline with less room left to move. Along the
\emph{syntactic-form} axis (novelty held fixed: Simile and Novel Metaphor are both
freshly-coined, compositional expressions with no prior conventional reading, differing
only in the explicit \emph{like}/\emph{as} marker), the contrast is sharp: Simile shows
the single largest effect in the entire study ($+0.256$), while Novel Metaphor shows none
on the same side. Holding novelty constant, syntactic packaging alone accounts for a
bigger swing than the entire conventionality axis does: an explicitly-marked
comparison transfers far more readily than a bare metaphorical assertion of identical
novelty. Taken together, the idiom vector's reach is more specific than a single
figurative-vs-literal axis in the abstract: it responds most to conventionalized
figurative senses and to explicitly-marked comparisons, and least to bare, freshly-coined
metaphor.

\textbf{Limitations.} (i) Our idiom benchmark is a hand-reconstructed 20-idiom set, not
IdioSteer's original stimuli, so absolute rates are not directly comparable across
papers, even though both benchmarks are now matched in size and both passed a
dual-annotator ambiguity-validation pass. (ii) Our judge-vs-human calibration covers only
86 idiom items and none of the figurative benchmark, well short of IdioSteer's own
280-item human gold set; a larger sample would give a more precise agreement number and
would let us check whether the same literal-leaning bias shows up on the figurative
benchmark too. (iii) The 10/5/5-expression category split is unequal and was not chosen
to match the conventionality/syntactic-form axes it is later read against (Simile and
Novel Metaphor happen to match at 5 each; Conventional Metaphor does not). Combined with
the resulting small per-category sample sizes (5 expressions $\times$ 5 variants for
Novel Metaphor and Simile), category-level deltas should be read as suggestive trends,
not precise, comparably-powered effect-size estimates. (iv) The Conventional Metaphor /
idiom boundary check (Section~\ref{sec:fig-data}) is a single-author substitutability
judgment, not independently cross-checked by a second annotator the way bidirectional
ambiguity is.

\section{Conclusion}
\label{sec:conclusion}

We asked whether the steering direction discovered by IdioSteer for idioms reflects a
general figurative-vs-literal axis, or is specific to idioms. Applying the idiom-derived
vector, unmodified, to a new, paper-scale (100-item), dual-annotator-validated benchmark
of non-idiomatic figurative expressions, we find broad, largely bidirectional
generalization, and the benchmark's two deliberately-varied axes (Section~\ref{sec:fig-data})
each explain part of it: along conventionality (form held fixed), transfer strength tracks
conventionality directly (Idiom $>$ Conventional Metaphor $>$ Novel Metaphor); along
syntactic form (novelty held fixed between Simile and Novel Metaphor), an explicitly-marked
comparison transfers far more readily than a bare metaphorical assertion of identical
novelty (Simile's $+0.256$, the largest effect in the study, vs. Novel Metaphor's
$\approx 0$ on the same side). So the idiom vector is not simply ``generalizes'' or
``doesn't'': it reaches non-idiomatic figurative language most when the expression is
either conventionalized or explicitly marked as a comparison, and least when it is both
novel and bare, plausibly a ceiling effect rather than true insensitivity given Novel
Metaphor's already literal-skewed baseline and rising incoherence at extreme coefficients.
Our own idiom replication is a good qualitative match to IdioSteer's reported
effect in shape, though its absolute baseline literal rate remains well above the
paper's own reported value, a gap our small-scale judge-vs-human calibration
($\kappa=0.494$ vs.\ IdioSteer's $\kappa=0.821$, with a literal-leaning judge bias)
suggests is partly attributable to judge quality, alongside a possibly less
figurative-skewed default register in our validated prefixes; we cannot fully separate
these two factors with the data collected so far. A promising next step, time permitting, is learning a
second steering vector directly from the figurative benchmark and comparing its geometry
(cosine similarity) to the idiom vector. If the two are highly aligned despite their
different category-level transfer strengths, that would be stronger evidence for a
shared underlying axis than the behavioral comparison alone can provide.

\bibliographystyle{plainnat}
\begin{thebibliography}{9}

\bibitem[Anonymous, 2026]{idiosteer}
Anonymous.
\emph{Steering Idioms: Controlling Figurative-vs-Literal Interpretation via
Residual-Stream Activation Steering}.
Anonymous ACL submission (under double-blind review), 2026. Provided as course
material.

\bibitem[Hashiloni et al., 2026]{idiolink}
Hashiloni et al.
\emph{IdioLink: Retrieving Meaning Beyond Words Across Idiomatic and Literal
Expressions}.
arXiv:2605.22247, 2026.

\bibitem[Turner et al., 2024]{actadd}
Turner et al.
\emph{Steering Language Models With Activation Addition}, 2024.

\bibitem[Rimsky et al., 2024]{caa}
Rimsky et al.
\emph{Steering Llama 2 via Contrastive Activation Addition}, 2024.

\bibitem[Haagsma et al., 2020]{magpie}
Haagsma et al.
\emph{MAGPIE: A Large Corpus of Potentially Idiomatic Expressions}, 2020.

\bibitem[Nunberg et al., 1994]{nunberg}
Nunberg, G., Sag, I. A., and Wasow, T.
\emph{Idioms}.
Language, 70(3):491--538, 1994.

\bibitem[Anonymous, 2026]{steerability}
Anonymous.
\emph{On the Steerability and Non-Surjectivity of Activation Steering
Directions}. ICLR 2026 submission. \url{https://iclr.cc/virtual/2026/10016066}

\end{thebibliography}

\end{document}
